\section[6. Numpy and Pandas]{6. Numpy and Pandas}

\subsection{Numpy}

\begin{itemize}
\item \textbf{스칼라}\textit{\textsuperscript{Scalar}}: 크기만 가지고 방향 가지지 않는 수. 실수 정수 복소수 등이 예시.
  한 개의 값
\item \textbf{벡터}\textit{\textsuperscript{Vector}}: 크기와 방향을 가지는 수. 여러 개의 값을 가지며 순서가 중요함.
\item \textbf{행렬}\textit{\textsuperscript{Matrix}}: 2차원 배열. 행과 열의 요소를 가짐. 각 요소는 스칼라이며, 행과 열에 따라 위치가 결정됨. 선형 변환, 선형 방정식, 데이터 표현 등에 널리 이용.
\end{itemize}

벡터나 행렬을 파이썬에서 다루기 위해서 주로 numpy를 많이 이용한다.\\
numpy는 파이썬에서 수치 계산을 위한 핵심 라이브러리로 널리 사용되는 도구로, 벡터와 행렬을 다루는데 굉장히 유용한 기능들이 다수 포함되어 있다. 즉, numpy를 사용하면 과학적이고 수치적인 계산을 효율적으로 수행할 수 있다.\\
우선, 기존에 벡터나 행렬을 표현하기 위해 사용했던 파이썬 리스트는 아래와 같은 특징이 있다.
\begin{itemize}
\item \textbf{높은 유연성}: list는 다양한 데이터 유형을 포함할 수 있으며 크기가 동적으로 조정될 수 있음.
\item \textbf{여러 내장 함수}: 다양한 내장 함수와 method가 있어 데이터 조작과 분석이 쉬움
\item \textbf{느린 속도}: 연산이 반복문 통해서 요소 별로 수행되기에 큰 데이터 셋에서 속도가 느림
\item \textbf{많은 메모리 사용}: 각 요소에 대한 정보 저장하기 위해 추가적인 메모리 공간을 이용하여 큰 데이터 셋에서 메모리 사용이 비효율적임
\end{itemize}

파이썬 리스트는 굉장히 느린 편이다. 이는 리스트가 내부 구현상으로는 연속된 나열로 된 값이 아니라 연속된 주소의 배열을 갖고, 각 주소가 리스트의 각기 다른 요소를 가리키는 방식이기 때문이다.\\
반면, numpy 배열은 파이썬 리스트와 비교하여 아래와 같은 특징이 있다.
\begin{itemize}
\item \textbf{빠른 성능}: 배열이 C로 구현되어 있기 때문에 반복문 쓰지 않고 벡터화된 연산을 수행할 수 있어 처리 속도가 빠름.
\item \textbf{적은 메모리 사용}: 연속된 공간에 데이터 저장하여 메모리 사용이 효율적
\item \textbf{다차원 배열의 지원}: 다차원 배열 지원하여 행렬 연산 및 다차원 데이터 처리가 용이
\item \textbf{다양한 수학 함수}: 다양한 수학 함수와 연산을 제공함
\item \textbf{낮은 유연성}: 한 가지 데이터 유형만 포함할 수 있고, 크기를 동적으로 조정할 수 없음. 크기 변경하려면 새로운 배열 만들어야 함.
\item \textbf{추가 기능의 부족}: 기본적인 수학과 배열 조작에 중점을 둔 함수밖에 없어서 추가 기능이 제한적일 수 있음.
\end{itemize}

numpy 배열은 파이썬 리스트에 비해서는 월등한 속도를 보여준다. 이때, 파이썬 리스트와 numpy 배열은 서로 장단점이 명확하게 구분되기 때문에, 시기적절하게 사용하는 것이 중요하다. 이때, numpy의 주요 함수는 다음과 같다.

\begin{longtable}{p{4cm}p{10cm}}
\toprule
함수 이름 & 설명\\
\midrule
array() & 배열 생성. 리스트나 튜플 객체 입력 받아 numpy 배열로 변환 \\
zeros() & 모든 요소가 0 \\
ones() & 모든 요소가 1 \\
arange() & 주어진 범위 내에서 일정한 간격으로 값이 증가하는 1차원 배열 \\
linspace() & 주어진 범위 내에서 지정된 개수에 맞춰서 일정한 간격으로 값 증가하는 1차원 배열 \\
reshape() & 배열의 차원과 크기 재조정 \\
transpose() & 전치 \\
concatenate() & 여러 배열을 연결함. 특정 축을 따라 배열을 연결할 수 있음. \\
sum() & 모든 요소의 합 \\
mean() & 모든 요소의 평균 \\
max, min() & 최대, 최소 \\
matmul() & 행렬 곱 \\
sort() & 배열 정렬. 특정 축 기준으로 할 수 있음. \\
argmax() & 최댓값의 index. 특정 축 기준으로 할 수 있음. \\
argmin() & 최솟값의 index. 특정 축 기준으로 할 수 있음. \\
\bottomrule
\end{longtable}

\subsection{Pandas}

pandas는 대규모 데이터를 조작하고 분석하기 위한 라이브러리이다. 강력한 데이터 구조와 데이터 조작 기능 제공하여 데이터를 쉽게 조작하고 분석할 수 있게 해준다.\\
\textbf{시리즈}\textit{\textsuperscript{Series}}는 1차원 배열과 같은 자료구조이다. 각 요소는 인덱스-값 쌍으로 구성되어 있으며 인덱스는 각 요소에 대해 고유한 숫자이다.\\
여러 개의 시리즈들를 모아서 \textbf{
데이터프레임}\textit{\textsuperscript{DataFrame}}을 구성할 수 있고, 이때 각 열은 서로 다른 데이터 유형을 가질 수 있다. 데이터프레임에서의 각 행은 하나의 데이터 쌍을 의미하고 행마다 인덱스를 가지고 있다.

\begin{longtable}{p{4cm}p{10cm}}
\toprule
함수 이름 & 설명\\
\midrule
pd.DataFrame() & 데이터프레임 생성 함수\\
head(n) & 첫 n개 행을 반환 (기본값은 5) \\
tails(n) & 첫 n개 열을 반환 (기본값은 5) \\
info() & 기본 정보를 출력(열의 데이터 타입, 결측치 유무, 메모리 사용량) \\
describe() & 숫자형 열에 대한 기술 통계 정보 출력(평균, 표준편차, 최대최소, 사분위수) \\
shape & 행, 열의 크기 반환 \\
columns & 열의 index 반환 \\
index & 행의 index 반환 \\
values & 값을 2차원 배열 형태로 반환 \\
loc{[}rl, cl{]} & label 이용하여 행과 열을 선택 \\
iloc{[}ri, ci{]} & 위치 이용하여 행과 열을 선택 \\
isnull() & 결측치 있는 위치를 확인하여 boolean 형태로 반환 \\
dropna() & 결측치 있는 행 또는 열을 삭제 \\
fillna(v) & 결측치를 value로 채움 \\
sort\_values(c) & 열 기준으로 데이터프레임 정렬. 기본적으로 오름차순이며, ascending=False로 내림차순으로 정렬 가능. \\
\bottomrule
\end{longtable}
